\documentclass[11pt]{article}
\usepackage[utf8]{inputenc}
\usepackage[T1]{fontenc}
\usepackage{amsmath,amssymb}
\usepackage{geometry}
\usepackage{listings}
\usepackage{xcolor}
\usepackage{hyperref}
\usepackage{tikz}
\usepackage{booktabs}

\geometry{margin=2.5cm}

\lstdefinelanguage{Rust}{
  keywords={fn,let,mut,for,in,if,else,while,loop,match,return,struct,impl,pub,use,mod,as,where,move,ref,self,Self,crate,super,async,await,unsafe,const,static,type,enum,trait, dyn, true, false},
  keywordstyle=\color{blue!70!black}\bfseries,
  ndkeywords={Vec,Vec3,Vec4,Quat,Mat4,u32,usize,f32,f64,bool,Option,Some,None,Result,Ok,Err,String},
  ndkeywordstyle=\color{purple!70!black}\bfseries,
  sensitive=true,
  comment=[l]{//},
  morecomment=[s]{/*}{*/},
  commentstyle=\color{green!50!black}\itshape,
  stringstyle=\color{red!70!black},
  morestring=[b]",
}

\lstset{
  basicstyle=\ttfamily\footnotesize,
  breaklines=true,
  frame=single,
  language=Rust,
  showstringspaces=false,
  tabsize=2,
}

\title{Orbit Camera: Euler-to-Quaternion Refactor}
\author{engvis}
\date{\today}

\begin{document}
\maketitle

\section{Problem Statement}

The orbit camera in \texttt{crates/engvis-core/src/camera.rs}
originally stored its orientation as two Euler angles:
\begin{itemize}
  \item \texttt{yaw} --- azimuth around the world $Y$ axis,
  \item \texttt{pitch} --- elevation, clamped to
        $\pm(\tfrac{\pi}{2} - 0.01) \approx \pm 89.4^\circ$.
\end{itemize}

When the user drags the mouse to rotate the view vertically, the
camera stops responding once \texttt{pitch} reaches the clamp
boundary.  The view cannot be rotated past nearly straight up/down,
which is perceived as a hard boundary on interaction.

\section{Root Cause Analysis}

\subsection{Pitch Clamp}

The clamp is required by the Euler representation itself, not by any
physical constraint of the scene.  In \texttt{orbit()}:

\begin{lstlisting}
pub fn orbit(&mut self, delta_yaw: f32, delta_pitch: f32) {
    self.yaw += delta_yaw;
    self.pitch = (self.pitch + delta_pitch).clamp(
        -std::f32::consts::FRAC_PI_2 + 0.01,
         std::f32::consts::FRAC_PI_2 - 0.01,
    );
}
\end{lstlisting}

Without the clamp, \texttt{pitch} would exceed $\pm\tfrac{\pi}{2}$,
which causes two failure modes:
\begin{enumerate}
  \item \textbf{Gimbal lock.}  At $\texttt{pitch} = \pm\tfrac{\pi}{2}$,
        the yaw and pitch rotation axes coincide, and the camera loses
        one degree of freedom.  Subsequent yaw input produces no
        visible rotation.
  \item \textbf{Up-vector flip.}  \texttt{view\_matrix()} uses
        \texttt{Mat4::look\_at\_rh(pos, target, Vec3::Y)}, which
        degenerates when the view direction aligns with $\pm Y$:
        the cross product $\text{forward} \times \text{up}$ vanishes,
        producing a NaN view matrix.
\end{enumerate}

\subsection{Why Euler Angles Are Inappropriate}

Euler angles parameterise $\mathrm{SO}(3)$ with a triple
$(\alpha, \beta, \gamma)$, but the map
\[
  (\alpha, \beta, \gamma) \mapsto R_z(\alpha) R_y(\beta) R_z(\gamma)
\]
is singular at $\beta = \pm \tfrac{\pi}{2}$ (gimbal lock).  Any
clamping strategy therefore imposes an artificial boundary on user
input, even though the underlying rotation space
$\mathrm{SO}(3) \cong \mathbb{RP}^3$ has no such boundary.

\section{Solution: Quaternion Orientation}

\subsection{Representation}

Replace the two Euler angles with a single unit quaternion
\texttt{orientation: Quat} that maps camera-local axes to world-space
directions:
\begin{itemize}
  \item camera $+X$ (right) $\mapsto$ \texttt{orientation * Vec3::X},
  \item camera $+Y$ (up)    $\mapsto$ \texttt{orientation * Vec3::Y},
  \item camera $+Z$ (back, away from target) $\mapsto$
        \texttt{orientation * Vec3::Z}.
\end{itemize}

The camera position is then
\[
  \mathbf{p}_{\text{cam}} = \mathbf{t} + \mathbf{q} \, (d \cdot \hat{\mathbf{z}}),
\]
where $\mathbf{t}$ is the orbit target, $\mathbf{q}$ the orientation
quaternion, and $d$ the orbit distance.

\begin{lstlisting}
pub struct OrbitCamera {
    pub target: Vec3,
    pub orientation: Quat,
    pub distance: f32,
    pub fov_y: f32,
    pub near: f32,
    pub far: f32,
    pub aspect_ratio: f32,
}

pub fn position(&self) -> Vec3 {
    self.target + self.orientation * (Vec3::Z * self.distance)
}
\end{lstlisting}

\subsection{Incremental Orbit Rotation}

Mouse drag produces two incremental rotations per frame:
\begin{itemize}
  \item \textbf{Yaw} $\Delta\alpha$ --- rotation about the world up
        axis $Y$.  Applied as a \emph{pre-multiplication}
        $\mathbf{q}' = R_Y(\Delta\alpha) \, \mathbf{q}$ so the horizon
        stays level regardless of the camera's current pitch.
  \item \textbf{Pitch} $\Delta\beta$ --- rotation about the camera's
        local $X$ axis.  Applied as a \emph{post-multiplication}
        $\mathbf{q}' \leftarrow \mathbf{q}' \, R_X(-\Delta\beta)$,
        where the sign preserves the old ``drag up = look up''
        convention.
\end{itemize}

\[
  \mathbf{q}_{\text{new}}
  = \mathrm{normalize}\!\left(
      R_Y(\Delta\alpha) \;\cdot\; \mathbf{q}_{\text{old}} \;\cdot\; R_X(-\Delta\beta)
    \right).
\]

Because quaternion multiplication is a closed, smooth operation on
$\mathrm{SO}(3)$, no clamp is required and no singularity arises at
any orientation.

\begin{lstlisting}
pub fn orbit(&mut self, delta_yaw: f32, delta_pitch: f32) {
    let yaw_rot   = Quat::from_rotation_y(delta_yaw);
    let pitch_rot = Quat::from_rotation_x(-delta_pitch);
    self.orientation = (yaw_rot * self.orientation * pitch_rot).normalize();
}
\end{lstlisting}

The final \texttt{.normalize()} guards against floating-point drift
after many incremental updates; it is a numerical safeguard, not a
behavioural clamp.

\subsection{View Matrix from Quaternion}

The previous implementation used
\texttt{Mat4::look\_at\_rh(position(), target, Vec3::Y)}, which
degenerates when $\text{forward} \parallel \text{up}$.  The new
implementation builds the view matrix directly from the orientation
quaternion:
\[
  V = R(\mathbf{q}^{-1}) \, T(-\mathbf{p}_{\text{cam}}),
\]
where $R(\mathbf{q}^{-1})$ is the world-to-camera rotation and
$T(-\mathbf{p}_{\text{cam}})$ translates the camera to the origin.

\begin{lstlisting}
pub fn view_matrix(&self) -> Mat4 {
    let rot  = Mat4::from_quat(self.orientation.inverse());
    let trans = Mat4::from_translation(-self.position());
    rot * trans
}
\end{lstlisting}

This formulation is well-defined for every orientation, including
straight up/down and upside-down views.

\subsection{Preset Views}

Each preset constructs the orientation quaternion directly from the
relevant axis rotations:

\begin{center}
\begin{tabular}{lll}
\toprule
Preset & Orientation & Description \\
\midrule
\texttt{view\_front} & $\mathbf{q} = I$ & camera at $+Z$, looking at $-Z$ \\
\texttt{view\_top}   & $\mathbf{q} = R_X(-\tfrac{\pi}{2})$ & camera at $+Y$, looking at $-Y$ \\
\texttt{view\_right} & $\mathbf{q} = R_Y(-\tfrac{\pi}{2})$ & camera at $-X$, looking at $+X$ \\
\texttt{view\_iso}   & $\mathbf{q} = R_Y(-0.785) \, R_X(-0.615)$ & isometric \\
\bottomrule
\end{tabular}
\end{center}

\begin{lstlisting}
pub fn view_front(&mut self) { self.orientation = Quat::IDENTITY; }
pub fn view_top(&mut self)   { self.orientation = Quat::from_rotation_x(-std::f32::consts::FRAC_PI_2); }
pub fn view_right(&mut self) { self.orientation = Quat::from_rotation_y(-std::f32::consts::FRAC_PI_2); }
pub fn view_iso(&mut self)  {
    self.orientation = Quat::from_rotation_y(-0.785) * Quat::from_rotation_x(-0.615);
}
\end{lstlisting}

\subsection{Compatibility Helper}

Code that previously set \texttt{yaw}/\texttt{pitch} directly (e.g.\
the \texttt{colorbar3d} example's reset-view button) now uses
\texttt{set\_orientation\_yaw\_pitch}, which composes the same
$R_Y \cdot R_X$ rotation:

\begin{lstlisting}
pub fn set_orientation_yaw_pitch(&mut self, yaw: f32, pitch: f32) {
    self.orientation = Quat::from_rotation_y(yaw) * Quat::from_rotation_x(-pitch);
}
\end{lstlisting}

\section{Comparison}

\begin{center}
\begin{tabular}{lll}
\toprule
Aspect & Euler (yaw, pitch) & Quaternion \\
\midrule
Representation        & $(\alpha, \beta)$, $\beta \in (-\tfrac{\pi}{2}, \tfrac{\pi}{2})$ & unit $\mathbf{q} \in S^3$ \\
Singularity            & at $\beta = \pm \tfrac{\pi}{2}$ (gimbal lock) & none \\
Pitch clamp           & required & not required \\
Vertical rotation     & stops at $\pm 89.4^\circ$ & unlimited \\
View matrix           & \texttt{look\_at\_rh} (degenerate at poles) & direct from $\mathbf{q}^{-1}$ \\
Up vector             & implicit ($+Y$) & explicit (\texttt{orientation * Vec3::Y}) \\
Preset views          & set two scalars & construct quaternion \\
Interpolation         & Euler interpolation flawed & SLERP / NLERP natural \\
Memory                & 8 bytes & 16 bytes \\
\bottomrule
\end{tabular}
\end{center}

\section{Pipeline Summary}

The complete camera update pipeline per frame:

\begin{enumerate}
  \item \textbf{Input} (\texttt{input.rs}) --- mouse drag deltas
        $(\Delta x, \Delta y)$ scaled by a sensitivity factor.
  \item \textbf{Orbit} --- \texttt{camera.orbit($-\Delta x \cdot s$,
        $-\Delta y \cdot s$)} composes incremental quaternion
        rotations.
  \item \item \textbf{Pan} (right mouse) --- translates
        \texttt{target} along camera-local $X$ and $Y$.
  \item \textbf{Zoom} (scroll) --- scales \texttt{distance} with
        clamp $[0.1, 500]$.
  \item \textbf{View matrix} --- built from
        $\mathbf{q}^{-1}$ and $-\mathbf{p}_{\text{cam}}$.
  \item \textbf{Projection} --- perspective RH from \texttt{fov\_y},
        \texttt{aspect\_ratio}, \texttt{near}, \texttt{far}.
\end{enumerate}

\section{Result}

After the refactor:
\begin{itemize}
  \item The camera rotates freely in all directions; no boundary is
        encountered when looking straight up or down.
  \item All preset views (front, top, right, iso) reproduce the
        previous camera framing.
  \item The \texttt{colorbar3d} example's reset-view button restores
        the same isometric orientation as before.
  \item \texttt{cargo build}, \texttt{cargo test --workspace}, and
        runtime startup all pass without errors.
\end{itemize}

\end{document}
